% ── ادامه نمودار اعتماد به نهادها ۲۰۲۲ ──────────────────
\begin{figurebox}
\begin{figure}[H]
\centering
\begin{tikzpicture}
\begin{axis}[
  width=13cm, height=9cm,
  xbar, bar width=10pt,
  xlabel={\rl{درصد اعتماد}},
  xmin=0, xmax=80,
  symbolic y coords={%
    {\rl{احزاب سیاسی}},
    {\rl{پارلمان}},
    {\rl{دولت مرکزی}},
    {\rl{قوه‌ی قضاییه}},
    {\rl{رسانه‌ها}},
    {\rl{نیروهای امنیتی}},
    {\rl{مراجع دینی}},
    {\rl{حشد الشعبی}},
    {\rl{ارتش}}},
  ytick=data,
  yticklabel style={font=\small},
  title={\rl{\textbf{اعتماد به نهادها (\lr{Arab Barometer} ۲۰۲۲)}}},
  title style={font=\small},
  nodes near coords,
  nodes near coords style={font=\small, anchor=west},
  enlarge y limits=0.1,
  % رنگ‌بندی بر اساس سطح اعتماد
  point meta=explicit symbolic,
  visualization depends on=\thisrow{color}\as\mycolor,
  scatter/@pre marker code/.append style={/tikz/fill=\mycolor},
]
\addplot[fill=red!45] coordinates {
  (8,{\rl{احزاب سیاسی}})
  (18,{\rl{پارلمان}})
  (22,{\rl{دولت مرکزی}})
};
\addplot[fill=yellow!55] coordinates {
  (25,{\rl{قوه‌ی قضاییه}})
  (35,{\rl{رسانه‌ها}})
};
\addplot[fill=blue!40] coordinates {
  (48,{\rl{نیروهای امنیتی}})
  (55,{\rl{مراجع دینی}})
  (58,{\rl{حشد الشعبی}})
};
\addplot[fill=green!50] coordinates {
  (72,{\rl{ارتش}})
};
\end{axis}
\end{tikzpicture}
\caption{ارتش (۷۲٪) و مراجع دینی (۵۵٪) بالاترین؛ احزاب (۸٪) و پارلمان (۱۸٪) پایین‌ترین}
\label{fig:bar-trust-2022}
\end{figure}
\end{figurebox}


% ████████████████████████████████████████████████████████████
%  بخش ۷ ─ اقتصاد، خدمات و کیفیت زندگی
% ████████████████████████████████████████████████████████████
\section{اقتصاد، خدمات و کیفیت زندگی}

\noindent
از سیاست به زندگی روزمره می‌رسیم.
سه منبع ─ \lr{Gallup}، \lr{IRI} و \lr{Arab Barometer} ─
ابعاد مختلف تجربه‌ی اقتصادی عراقی‌ها را پوشش می‌دهند.


\subsection{نردبان زندگی کنتریل: عراق در مقایسه‌ی جهانی}

\begin{quote}\small
\lr{``Please imagine a ladder with steps from zero (worst possible
life) to ten (best possible life).
On which step do you stand now?
On which step will you stand five years from now?''}
\end{quote}

\begin{table}[H]
\centering
\caption{نردبان زندگی کنتریل ─ عراق در مقایسه (\lr{Gallup})}
\label{tab:gallup-cantril}
\small
\begin{tabularx}{\textwidth}{>{\bfseries\small}r *{6}{>{\centering\arraybackslash\small}X}}
\toprule
& \textbf{۲۰۰۶} & \textbf{۲۰۰۷} & \textbf{۲۰۰۸} & \textbf{۲۰۰۹} & \textbf{۲۰۱۰} & \textbf{۲۰۱۱} \\
\midrule
\multicolumn{7}{r}{\textbf{زندگی فعلی (میانگین از ۱۰)}} \\
عراق           & \pdec{۴}{۱} & \pdec{۳}{۸} & \pdec{۴}{۲} & \pdec{۴}{۵} & \pdec{۴}{۷} & \pdec{۴}{۶} \\
میانگین جهانی  & \pdec{۵}{۳} & \pdec{۵}{۳} & \pdec{۵}{۴} & \pdec{۵}{۳} & \pdec{۵}{۴} & \pdec{۵}{۴} \\
افغانستان      & \pdec{۳}{۸} & \pdec{۳}{۹} & \pdec{۴}{۰} & \pdec{۴}{۲} & \pdec{۴}{۳} & \pdec{۴}{۴} \\
\midrule
\multicolumn{7}{r}{\textbf{زندگی ۵ سال آینده (میانگین از ۱۰)}} \\
عراق           & \pdec{۵}{۵} & \pdec{۵}{۲} & \pdec{۵}{۸} & \pdec{۶}{۰} & \pdec{۶}{۱} & \pdec{۵}{۸} \\
میانگین جهانی  & \pdec{۶}{۶} & \pdec{۶}{۵} & \pdec{۶}{۶} & \pdec{۶}{۶} & \pdec{۶}{۶} & \pdec{۶}{۶} \\
\midrule
\multicolumn{7}{r}{\textbf{طبقه‌بندی گالوپ}} \\
شکوفا (\lr{Thriving})          & ۱۷٪ & ۱۲٪ & ۱۶٪ & ۲۲٪ & ۲۳٪ & ۲۲٪ \\
دست‌وپنجه (\lr{Struggling})    & ۶۲٪ & ۶۰٪ & ۶۰٪ & ۵۸٪ & ۵۸٪ & ۵۹٪ \\
رنج‌کشیده (\lr{Suffering})     & ۲۱٪ & ۲۸٪ & ۲۴٪ & ۲۰٪ & ۱۹٪ & ۱۹٪ \\
\bottomrule
\end{tabularx}
\end{table}

\begin{figurebox}
\begin{figure}[H]
\centering
\begin{tikzpicture}
\begin{axis}[
  width=14cm, height=7cm,
  xlabel={\rl{سال}},
  ylabel={\rl{میانگین (از ۱۰)}},
  xmin=2005.5, xmax=2011.5,
  ymin=3, ymax=7.5,
  xtick={2006,2007,2008,2009,2010,2011},
  legend pos=north east,
  legend style={font=\small},
  grid=major, grid style={dashed, gray!25},
  title={\rl{\textbf{نردبان زندگی کنتریل: عراق همیشه زیر میانگین جهانی}}},
  title style={font=\small},
]
\addplot[color=red, mark=*, very thick, mark size=3pt] coordinates {
  (2006,4.1)(2007,3.8)(2008,4.2)(2009,4.5)(2010,4.7)(2011,4.6)
};
\addlegendentry{\rl{عراق ─ فعلی}}
\addplot[color=red, mark=*, very thick, mark size=3pt, dashed] coordinates {
  (2006,5.5)(2007,5.2)(2008,5.8)(2009,6.0)(2010,6.1)(2011,5.8)
};
\addlegendentry{\rl{عراق ─ ۵ سال آینده}}
\addplot[color=gray, mark=square*, thick, mark size=2pt] coordinates {
  (2006,5.3)(2007,5.3)(2008,5.4)(2009,5.3)(2010,5.4)(2011,5.4)
};
\addlegendentry{\rl{میانگین جهانی ─ فعلی}}
\addplot[color=orange, mark=triangle*, thick, mark size=2pt, dotted] coordinates {
  (2006,3.8)(2007,3.9)(2008,4.0)(2009,4.2)(2010,4.3)(2011,4.4)
};
\addlegendentry{\rl{افغانستان ─ فعلی}}
\end{axis}
\end{tikzpicture}
\caption{عراق بین افغانستان و میانگین جهانی ─ اما خوش‌بینی به آینده بالاتر از واقعیت}
\label{fig:gallup-cantril}
\end{figure}
\end{figurebox}

\didyouknow{%
  فاصله‌ی بین «زندگی فعلی» و «۵ سال آینده» در عراق
  به‌طور مداوم بیش از ۱ واحد بود.
  این «شکاف امید» (\lr{Hope Gap}) نشان می‌دهد عراقی‌ها
  حتی در بدترین شرایط، باور به بهبود آینده را از دست
  ندادند ─ سرمایه‌ای اجتماعی که هنوز قابل بهره‌برداری است.
}


\subsection{احساس امنیت شخصی و تجربه‌ی خشونت}

\begin{comparecard}[title={\textbf{گالوپ: امنیت شخصی و تجربه‌ی خشونت}}]
% ── ستون چپ: امنیت شبانه ──
\begin{center}\textbf{«آیا شب‌ها احساس امنیت می‌کنید؟»}\end{center}
\small
\begin{tabularx}{0.95\linewidth}{>{\bfseries\scriptsize}r *{4}{>{\centering\arraybackslash\scriptsize}X}}
\toprule
& \textbf{۰۷} & \textbf{۰۹} & \textbf{۱۱} & \textbf{جهانی} \\
\midrule
بله (کل) & ۱۸٪ & ۴۲٪ & ۴۸٪ & ۶۳٪ \\
مرد      & ۲۲٪ & ۵۰٪ & ۵۵٪ & -- \\
زن       & ۱۴٪ & ۳۴٪ & ۴۱٪ & -- \\
\bottomrule
\end{tabularx}
\tcblower
% ── ستون راست: تجربه خشونت ──
\begin{center}\textbf{تجربه‌ی خشونت و محرومیت}\end{center}
\small
\begin{tabularx}{0.95\linewidth}{>{\bfseries\scriptsize}r *{3}{>{\centering\arraybackslash\scriptsize}X}}
\toprule
& \textbf{۰۷} & \textbf{۰۹} & \textbf{۱۱} \\
\midrule
قربانی حمله/سرقت   & ۳۲٪ & ۱۸٪ & ۱۳٪ \\
پول کافی غذا نداشتم & ۳۵٪ & ۲۵٪ & ۲۰٪ \\
عضو خانواده کشته شده & ۲۲٪ & ۲۰٪ & ۱۸٪ \\
مجبور به ترک خانه   & ۱۸٪ & ۱۵٪ & ۱۰٪ \\
\bottomrule
\end{tabularx}
\end{comparecard}


\subsection{ارزیابی اقتصادی: گالوپ}

\begin{table}[H]
\centering
\caption{ارزیابی اقتصادی ─ گالوپ}
\label{tab:gallup-economy}
\small
\begin{tabularx}{\textwidth}{>{\bfseries\small}r *{6}{>{\centering\arraybackslash\small}X}}
\toprule
& \textbf{۲۰۰۶} & \textbf{۲۰۰۷} & \textbf{۲۰۰۸} & \textbf{۲۰۰۹} & \textbf{۲۰۱۰} & \textbf{۲۰۱۱} \\
\midrule
اقتصاد محلی بهتر می‌شود        & ۲۲٪ & ۱۸٪ & ۲۸٪ & ۳۵٪ & ۳۸٪ & ۳۲٪ \\
وضعیت اقتصادی «عالی/خوب»      & ۱۲٪ & ۱۰٪ & ۱۵٪ & ۲۰٪ & ۲۲٪ & ۱۸٪ \\
رضایت از سطح زندگی             & ۵۵٪ & ۴۸٪ & ۵۲٪ & ۵۸٪ & ۶۰٪ & ۵۸٪ \\
بازار کار محلی خوب              & ۱۲٪ & ۱۰٪ & ۱۵٪ & ۱۸٪ & ۲۰٪ & ۱۸٪ \\
\midrule
\multicolumn{7}{r}{\textbf{بازار کار ─ مقایسه}} \\
اردن           & ۲۵٪ & ۲۲٪ & ۲۰٪ & ۱۸٪ & ۲۰٪ & ۲۲٪ \\
میانگین جهانی  & ۳۵٪ & ۳۲٪ & ۳۰٪ & ۲۸٪ & ۲۸٪ & ۳۰٪ \\
\bottomrule
\end{tabularx}
\end{table}


\subsection{شکست خدمات عمومی: \lr{IRI}}

\begin{surveycard}{\lr{International Republican Institute (IRI)} ─ ۲۰۰۳–۲۰۱۳}
\begin{tabularx}{\textwidth}{>{\bfseries\small}r >{\small}X}
هدف     & ارزیابی دموکراسی‌سازی و عملکرد دولت \\
نمونه   & \ptho{۳}{۰۰۰}–\ptho{۳}{۵۰۰} (بزرگ‌ترین نمونه‌ها) \\
موج‌ها   & ۱۲ موج \\
ویژگی   & پرسش‌های بسیار عملیاتی درباره‌ی خدمات دولتی \\
\end{tabularx}
\end{surveycard}

\begin{table}[H]
\centering
\caption{ارزیابی خدمات عمومی (درصد «خوب/عالی») ─ \lr{IRI}}
\label{tab:iri-services}
\small
\begin{tabularx}{\textwidth}{>{\bfseries\small}r *{6}{>{\centering\arraybackslash\small}X}}
\toprule
\textbf{خدمت} & \textbf{۲۰۰۴} & \textbf{۲۰۰۶} & \textbf{۲۰۰۸} & \textbf{۲۰۱۰} & \textbf{۲۰۱۱} & \textbf{۲۰۱۳} \\
\midrule
برق              & ۱۵٪ & ۱۲٪ & ۱۸٪ & ۲۰٪ & ۲۲٪ & ۲۰٪ \\
آب آشامیدنی      & ۳۵٪ & ۳۰٪ & ۳۵٪ & ۳۸٪ & ۴۰٪ & ۳۸٪ \\
فاضلاب           & ۲۰٪ & ۱۸٪ & ۲۲٪ & ۲۵٪ & ۲۵٪ & ۲۲٪ \\
بهداشت و درمان   & ۲۲٪ & ۱۸٪ & ۲۵٪ & ۲۸٪ & ۳۰٪ & ۲۵٪ \\
آموزش            & ۳۰٪ & ۲۵٪ & ۳۵٪ & ۳۸٪ & ۳۸٪ & ۳۵٪ \\
فرصت‌های شغلی    & ۱۰٪ & ۸٪  & ۱۲٪ & ۱۵٪ & ۱۵٪ & ۱۲٪ \\
\midrule
\multicolumn{7}{r}{\textbf{برق ─ تفکیک استانی (۲۰۱۱)}} \\
اربیل     & & & & & ۵۵٪ & \\
سلیمانیه  & & & & & ۵۰٪ & \\
نجف       & & & & & ۱۸٪ & \\
بغداد     & & & & & ۱۵٪ & \\
بصره      & & & & & ۱۲٪ & \\
انبار      & & & & & ۱۰٪ & \\
\bottomrule
\end{tabularx}
\end{table}

\infobox{boxred}{شکست خدمات عمومی پس از ۶۰ میلیارد دلار بازسازی}{%
حتی ده سال پس از ۲۰۰۳:
\begin{itemize}[rightmargin=0.5cm]
\item فقط \keynumber{red}{۲۰٪} برق را «خوب» ارزیابی کردند
\item فقط \keynumber{red}{۱۲٪} فرصت‌های شغلی را مناسب دانستند
\item تفاوت فاحش: کردستان (\keynumber{green}{۵۵٪} رضایت از برق)
  در برابر بقیه‌ی عراق (\keynumber{red}{۱۰–۱۸٪})
\item این ارقام پس از صرف بیش از ۶۰ میلیارد دلار بودجه‌ی بازسازی حاصل شد
\end{itemize}
}


\subsection{ارزیابی عملکرد دولت: بارومتر عرب}

\begin{table}[H]
\centering
\caption{ارزیابی عملکرد دولت (درصد «خوب/بسیار خوب») ─ بارومتر عرب}
\label{tab:ab-govt-performance}
\small
\begin{tabularx}{\textwidth}{>{\bfseries\small}r *{4}{>{\centering\arraybackslash\small}X}}
\toprule
\textbf{حوزه} & \textbf{III (۱۳)} & \textbf{IV (۱۶)} & \textbf{V (۱۹)} & \textbf{VII (۲۲)} \\
\midrule
ایجاد شغل            & ۸٪  & ۵٪  & ۵٪  & ۸٪  \\
کنترل قیمت‌ها         & ۱۲٪ & ۸٪  & ۸٪  & ۱۰٪ \\
مبارزه با فساد       & ۱۰٪ & ۵٪  & ۳٪  & ۸٪  \\
کاهش شکاف فقیر/غنی  & ۸٪  & ۵٪  & ۳٪  & ۵٪  \\
خدمات بهداشتی        & ۱۸٪ & ۱۲٪ & ۱۰٪ & ۱۵٪ \\
خدمات آموزشی         & ۲۰٪ & ۱۵٪ & ۱۲٪ & ۱۵٪ \\
تأمین برق            & ۱۵٪ & ۱۰٪ & ۱۲٪ & ۱۵٪ \\
تأمین آب             & ۱۸٪ & ۱۲٪ & ۱۵٪ & ۱۸٪ \\
\bottomrule
\end{tabularx}
\end{table}

\subsection{ارزیابی عملکرد مقامات منتخب: \lr{IRI}}

\begin{table}[H]
\centering
\caption{ارزیابی عملکرد مقامات (درصد «خوب/عالی») ─ \lr{IRI}}
\label{tab:iri-officials}
\small
\begin{tabularx}{\textwidth}{>{\bfseries\small}r *{6}{>{\centering\arraybackslash\small}X}}
\toprule
\textbf{مقام} & \textbf{۲۰۰۴} & \textbf{۲۰۰۶} & \textbf{۲۰۰۸} & \textbf{۲۰۱۰} & \textbf{۲۰۱۱} & \textbf{۲۰۱۳} \\
\midrule
نخست‌وزیر          & ۵۰٪ & ۴۲٪ & ۴۵٪ & ۴۸٪ & ۳۵٪ & ۲۸٪ \\
استاندار           & ۳۵٪ & ۳۰٪ & ۳۵٪ & ۳۸٪ & ۳۲٪ & ۲۸٪ \\
شورای استان        & ۲۸٪ & ۲۲٪ & ۲۸٪ & ۳۰٪ & ۲۵٪ & ۲۰٪ \\
شهردار/شورای محلی  & ۳۲٪ & ۲۵٪ & ۳۲٪ & ۳۵٪ & ۲۸٪ & ۲۵٪ \\
نماینده‌ی پارلمان   & ۲۵٪ & ۱۸٪ & ۲۲٪ & ۲۵٪ & ۱۸٪ & ۱۵٪ \\
\midrule
\multicolumn{7}{r}{\textbf{نخست‌وزیر ─ تفکیک قومی–مذهبی (۲۰۱۱: مالکی)}} \\
شیعه & & & & & ۴۵٪ & \\
سنی  & & & & & ۱۲٪ & \\
کرد  & & & & & ۳۵٪ & \\
\bottomrule
\end{tabularx}
\end{table}


% ████████████████████████████████████████████████████████████
%  بخش ۸ ─ اولویت‌های مردم: از امنیت تا فساد
% ████████████████████████████████████████████████████████████
\section{اولویت‌های مردم: چرخش بزرگ از امنیت به فساد}

\noindent
یکی از مهم‌ترین تحولات دو دهه‌ی گذشته، جابه‌جایی «دغدغه‌ی اول»
از امنیت به فساد است. دو منبع مستقل ─ \lr{ABC/BBC} و
\lr{Arab Barometer} ─ این چرخش را تأیید می‌کنند.

\subsection{مهم‌ترین مشکل: \lr{ABC/BBC}}

\begin{table}[H]
\centering
\caption{مهم‌ترین مشکل کشور (\lr{ABC/BBC})}
\label{tab:abc-q8-priorities}
\begin{tabularx}{\textwidth}{>{\bfseries\small}r *{5}{>{\centering\arraybackslash\small}X}}
\toprule
\textbf{مشکل} & \textbf{۲۰۰۴} & \textbf{۲۰۰۵} & \textbf{۲۰۰۷} & \textbf{۲۰۰۸} & \textbf{۲۰۰۹} \\
\midrule
امنیت/خشونت             & ۴۵٪ & ۵۰٪ & ۵۵٪ & ۴۵٪ & ۲۵٪ \\
بیکاری                  & ۲۰٪ & ۱۸٪ & ۱۲٪ & ۱۵٪ & ۲۲٪ \\
خدمات عمومی (برق، آب)   & ۱۵٪ & ۱۲٪ & ۸٪  & ۱۲٪ & ۱۸٪ \\
فساد                    & ۵٪  & ۵٪  & ۵٪  & ۱۰٪ & ۲۰٪ \\
اشغال خارجی             & ۸٪  & ۱۰٪ & ۱۵٪ & ۱۰٪ & ۵٪  \\
تروریسم                 & ۵٪  & ۳٪  & ۳٪  & ۵٪  & ۵٪  \\
تنش فرقه‌ای             & ۲٪  & ۲٪  & ۲٪  & ۳٪  & ۵٪  \\
\bottomrule
\end{tabularx}
\end{table}


\subsection{مهم‌ترین چالش: بارومتر عرب}

\begin{table}[H]
\centering
\caption{مهم‌ترین چالش ─ بارومتر عرب}
\label{tab:ab-challenges}
\begin{tabularx}{\textwidth}{>{\bfseries\small}r *{6}{>{\centering\arraybackslash\small}X}}
\toprule
\textbf{چالش} & \textbf{I} & \textbf{II} & \textbf{III} & \textbf{IV} & \textbf{V} & \textbf{VII} \\
\midrule
امنیت/تروریسم    & ۴۸٪ & ۲۵٪ & ۳۵٪ & ۴۰٪ & ۱۰٪ & ۸٪  \\
فساد مالی/اداری  & ۸٪  & ۲۵٪ & ۲۸٪ & ۲۵٪ & ۴۵٪ & ۳۸٪ \\
بیکاری           & ۱۵٪ & ۲۰٪ & ۱۵٪ & ۱۵٪ & ۲۰٪ & ۲۵٪ \\
وضعیت اقتصادی   & ۱۲٪ & ۱۵٪ & ۱۰٪ & ۱۰٪ & ۱۵٪ & ۱۸٪ \\
خدمات عمومی      & ۱۰٪ & ۸٪  & ۵٪  & ۵٪  & ۵٪  & ۵٪  \\
مداخله‌ی خارجی   & ۵٪  & ۵٪  & ۵٪  & ۳٪  & ۳٪  & ۴٪  \\
\bottomrule
\end{tabularx}
\end{table}


% ── نمودار بصری تغییر اولویت‌ها (جایگزین پای‌چارت) ──────
\begin{figurebox}
\begin{figure}[H]
\centering
\begin{tikzpicture}
\begin{axis}[
  width=15cm, height=8cm,
  ybar stacked,
  bar width=20pt,
  xlabel={\rl{دوره}},
  ylabel={\rl{درصد}},
  ymin=0, ymax=105,
  symbolic x coords={%
    {\rl{۲۰۰۶ (اوج جنگ)}},
    {\rl{۲۰۱۳ (پیش‌داعش)}},
    {\rl{۲۰۱۹ (تشرین)}},
    {\rl{۲۰۲۲ (ثبات نسبی)}}},
  xtick=data,
  xticklabel style={font=\small},
  legend pos=outer north east,
  legend style={font=\small, cells={anchor=west}},
  title={\rl{\textbf{تغییر اولویت‌ها: از امنیت (قرمز) به فساد (بنفش)}}},
  title style={font=\small},
  nodes near coords,
  nodes near coords style={font=\tiny},
]
\addplot[fill=red!55] coordinates {
  ({\rl{۲۰۰۶ (اوج جنگ)}},48)
  ({\rl{۲۰۱۳ (پیش‌داعش)}},35)
  ({\rl{۲۰۱۹ (تشرین)}},10)
  ({\rl{۲۰۲۲ (ثبات نسبی)}},8)
};
\addlegendentry{\rl{امنیت/تروریسم}}
\addplot[fill=purple!55] coordinates {
  ({\rl{۲۰۰۶ (اوج جنگ)}},8)
  ({\rl{۲۰۱۳ (پیش‌داعش)}},28)
  ({\rl{۲۰۱۹ (تشرین)}},45)
  ({\rl{۲۰۲۲ (ثبات نسبی)}},38)
};
\addlegendentry{\rl{فساد}}
\addplot[fill=orange!45] coordinates {
  ({\rl{۲۰۰۶ (اوج جنگ)}},15)
  ({\rl{۲۰۱۳ (پیش‌داعش)}},15)
  ({\rl{۲۰۱۹ (تشرین)}},20)
  ({\rl{۲۰۲۲ (ثبات نسبی)}},25)
};
\addlegendentry{\rl{بیکاری}}
\addplot[fill=gray!35] coordinates {
  ({\rl{۲۰۰۶ (اوج جنگ)}},29)
  ({\rl{۲۰۱۳ (پیش‌داعش)}},22)
  ({\rl{۲۰۱۹ (تشرین)}},25)
  ({\rl{۲۰۲۲ (ثبات نسبی)}},29)
};
\addlegendentry{\rl{سایر}}
\end{axis}
\end{tikzpicture}
\caption{چرخش اولویت‌ها: امنیت از ۴۸٪ به ۸٪ · فساد از ۸٪ به ۴۵٪}
\label{fig:stacked-priorities}
\end{figure}
\end{figurebox}

\infobox{boxblue}{تحلیل: چرخش از امنیت به فساد}{%
\begin{itemize}[rightmargin=0.5cm]
\item با کاهش نسبی خشونت، شهروندان فرصت یافتند به مشکلات ساختاری بیندیشند
\item فساد سیستماتیک نه‌تنها کاهش نیافته بلکه نهادینه‌تر شده
\item جنبش تشرین ۲۰۱۹ بر همین بستر نارضایتی از فساد شکل گرفت
\item مردم دریافته‌اند بسیاری از مشکلات امنیتی خود ریشه در فساد دارد
\end{itemize}
}


% ████████████████████████████████████████████████████████████
%  بخش ۹ ─ دموکراسی: ایده‌ی ماندگار، عملکرد ناامیدکننده
% ████████████████████████████████████████████████████████████
\section{دموکراسی: ایده‌ی ماندگار، عملکرد ناامیدکننده}

\noindent
شاید شگفت‌انگیزترین یافته‌ی نظرسنجی‌ها آن باشد که
عراقی‌ها بین \textbf{ایده‌ی دموکراسی} و \textbf{عملکرد آن}
تمایز قائل شده‌اند: از دومی سرخورده‌اند اما از اولی دست نکشیده‌اند.


\subsection{دموکراسی به‌مثابه نظام مطلوب: بارومتر عرب}

\begin{quote}\small
\lr{``A democratic system may have problems, yet it is better
than other systems. To what extent do you agree?''}
\end{quote}

\begin{table}[H]
\centering
\caption{«دموکراسی بهتر از سایر نظام‌هاست» ─ بارومتر عرب}
\label{tab:ab-democracy-best}
\small
\begin{tabularx}{\textwidth}{>{\bfseries\small}r *{6}{>{\centering\arraybackslash\small}X}}
\toprule
& \textbf{I (۰۶)} & \textbf{II (۱۱)} & \textbf{III (۱۳)} & \textbf{IV (۱۶)} & \textbf{V (۱۹)} & \textbf{VII (۲۲)} \\
\midrule
\multicolumn{7}{r}{\textbf{کل ─ مجموع موافق}} \\
مجموع موافق & ۵۹٪ & ۵۳٪ & ۵۳٪ & ۴۸٪ & ۵۵٪ & ۵۰٪ \\
\midrule
\multicolumn{7}{r}{\textbf{تفکیک سنی}} \\
۱۸–۲۹ & ۵۵٪ & ۵۰٪ & ۵۰٪ & ۴۵٪ & ۵۲٪ & ۴۸٪ \\
۳۰–۴۹ & ۶۰٪ & ۵۵٪ & ۵۵٪ & ۵۰٪ & ۵۸٪ & ۵۲٪ \\
۵۰+   & ۶۲٪ & ۵۵٪ & ۵۳٪ & ۵۰٪ & ۵۵٪ & ۵۱٪ \\
\midrule
\multicolumn{7}{r}{\textbf{تفکیک تحصیلاتی}} \\
ابتدایی و کمتر & ۵۰٪ & ۴۵٪ & ۴۵٪ & ۴۰٪ & ۴۵٪ & ۴۲٪ \\
متوسطه        & ۵۸٪ & ۵۲٪ & ۵۲٪ & ۴۸٪ & ۵۵٪ & ۵۰٪ \\
دانشگاهی      & ۶۸٪ & ۶۲٪ & ۶۰٪ & ۵۵٪ & ۶۵٪ & ۵۸٪ \\
\midrule
\multicolumn{7}{r}{\textbf{مقایسه‌ی منطقه‌ای}} \\
اردن   & ۸۰٪ & ۷۵٪ & ۷۰٪ & ۶۵٪ & ۶۰٪ & ۵۸٪ \\
لبنان  & ۸۵٪ & ۸۰٪ & ۷۵٪ & ۷۰٪ & ۶۸٪ & ۶۲٪ \\
مصر    & ۷۲٪ & ۷۰٪ & ۶۵٪ & ۵۵٪ & ۵۲٪ & ۵۰٪ \\
\bottomrule
\end{tabularx}
\end{table}


\subsection{آیا دموکراسی در عراق کار می‌کند؟ \lr{Pew}}

\begin{surveycard}{\lr{Pew Research Center ─ Global Attitudes} ─ ۲۰۰۳–۲۰۱۵}
\begin{tabularx}{\textwidth}{>{\bfseries\small}r >{\small}X}
مرکز   & \lr{Pew Research Center} (واشنگتن) \\
نمونه  & \ptho{۱}{۰۰۰}–\ptho{۱}{۵۰۰} \\
سال‌ها  & ۲۰۰۳–۰۹ · ۲۰۱۲ · ۲۰۱۵ \\
\end{tabularx}
\end{surveycard}

\begin{quote}\small
\lr{``Some believe democracy is a Western way that won't work here.
Others feel it can work. Which is closer to your view?''}
\end{quote}

\begin{table}[H]
\centering
\caption{«آیا دموکراسی در اینجا کار می‌کند؟» ─ \lr{Pew}}
\label{tab:pew-democracy-works}
\begin{tabularx}{\textwidth}{>{\bfseries}r *{5}{>{\centering\arraybackslash}X}}
\toprule
& \textbf{۲۰۰۴} & \textbf{۲۰۰۵} & \textbf{۲۰۰۷} & \textbf{۲۰۰۹} & \textbf{۲۰۱۲} \\
\midrule
کار می‌کند     & ۶۰٪ & ۵۵٪ & ۴۸٪ & ۵۵٪ & ۵۲٪ \\
کار نمی‌کند    & ۳۰٪ & ۳۵٪ & ۴۰٪ & ۳۵٪ & ۳۸٪ \\
\midrule
\multicolumn{6}{r}{\textbf{«کار می‌کند» ─ تفکیک قومی–مذهبی}} \\
کرد   & ۷۵٪ & ۷۰٪ & ۶۸٪ & ۷۰٪ & ۶۸٪ \\
شیعه  & ۶۵٪ & ۶۰٪ & ۵۰٪ & ۶۰٪ & ۵۵٪ \\
سنی   & ۴۰٪ & ۳۵٪ & ۳۰٪ & ۳۵٪ & ۳۵٪ \\
\bottomrule
\end{tabularx}
\end{table}


\subsection{رضایت از دموکراسی: روند نزولی}

\begin{figurebox}
\begin{figure}[H]
\centering
\begin{tikzpicture}
\begin{axis}[
  width=14cm, height=7cm,
  xlabel={\rl{سال}},
  ylabel={\rl{درصد}},
  xmin=2004, xmax=2023,
  ymin=0, ymax=85,
  xtick={2005,2007,2009,2011,2013,2015,2017,2019,2021,2023},
  xticklabel style={rotate=45, anchor=east, font=\small},
  legend pos=north east,
  legend style={font=\small, cells={anchor=west}},
  grid=major, grid style={dashed, gray!20},
  title={\rl{\textbf{رضایت از عملکرد دموکراسی در عراق}}},
  title style={font=\small},
]
\addplot[color=green!60!black, mark=*, thick, mark size=2pt] coordinates {
  (2005,55)(2007,30)(2009,35)(2011,38)
  (2013,28)(2015,15)(2017,20)(2019,12)(2021,18)(2022,22)
};
\addlegendentry{\rl{راضی/نسبتاً راضی}}
\addplot[color=red!70, mark=square*, thick, mark size=2pt] coordinates {
  (2005,30)(2007,55)(2009,50)(2011,48)
  (2013,60)(2015,72)(2017,65)(2019,78)(2021,70)(2022,65)
};
\addlegendentry{\rl{ناراضی/بسیار ناراضی}}
\draw[dashed, gray] (axis cs:2014,0) -- (axis cs:2014,80)
  node[above, font=\tiny] {\rl{داعش}};
\draw[dashed, orange] (axis cs:2019,0) -- (axis cs:2019,80)
  node[above, font=\tiny] {\rl{تشرین}};
\end{axis}
\end{tikzpicture}
\caption{رضایت از دموکراسی: از ۵۵٪ (۲۰۰۵) به ۲۲٪ (۲۰۲۲) ─ نارضایتی در تشرین ۷۸٪}
\label{fig:line-democracy-satisfaction}
\end{figure}
\end{figurebox}

\infobox{boxgreen}{پارادوکس دموکراسی عراقی}{%
علی‌رغم سرخوردگی عمیق از عملکرد:
\begin{itemize}[rightmargin=0.5cm]
\item حمایت از \textbf{اصل دموکراسی} هرگز به زیر ۴۸٪ نرفت
\item تفاوت بین «دموکراسی خوب است» (≈۵۰٪) و «راضی‌ام» (≈۲۲٪) = \textbf{شکاف ۲۸ واحدی}
\item این سرمایه‌ی اجتماعی نشان می‌دهد اعتمادسازی نهادی هنوز ممکن است
\end{itemize}
}


% ████████████████████████████████████████████████████████████
%  بخش ۱۰ ─ هویت، ارزش‌ها و همزیستی
% ████████████████████████████████████████████████████████████
\section{هویت، ارزش‌ها و همزیستی}

\noindent
از سیاست و اقتصاد به لایه‌های عمیق‌تر می‌رسیم:
هویت ملی در برابر فرقه‌ای، ارزش‌های فرهنگی و امکان همزیستی.


\subsection{هویت ملی در برابر فرقه‌ای}

\begin{quote}\small
\lr{``Which do you identify with most ─ being Iraqi, Muslim,
Shia/Sunni/Kurd, or Arab/Kurdish?''} (\lr{ABC/BBC})
\end{quote}

\begin{table}[H]
\centering
\caption{اولویت هویتی عراقی‌ها (\lr{ABC/BBC})}
\label{tab:abc-q9-identity}
\begin{tabularx}{\textwidth}{>{\bfseries}r *{3}{>{\centering\arraybackslash}X}}
\toprule
\textbf{هویت اول} & \textbf{۲۰۰۷} & \textbf{۲۰۰۸} & \textbf{۲۰۰۹} \\
\midrule
عراقی                    & ۲۸٪ & ۳۳٪ & ۴۱٪ \\
مسلمان                   & ۳۵٪ & ۳۰٪ & ۲۸٪ \\
مذهب (شیعه/سنی)          & ۲۲٪ & ۲۰٪ & ۱۵٪ \\
قومیت (عرب/کرد)          & ۱۵٪ & ۱۷٪ & ۱۶٪ \\
\midrule
\multicolumn{4}{r}{\textbf{هویت «عراقی» ─ تفکیک}} \\
سنی   & ۳۵٪ & ۴۲٪ & ۵۰٪ \\
شیعه  & ۳۰٪ & ۳۵٪ & ۴۵٪ \\
کرد   & ۸٪  & ۱۰٪ & ۱۲٪ \\
\bottomrule
\end{tabularx}
\end{table}


% ── هویت: میله‌ای انباشته (جایگزین پای‌چارت ۵) ──────────
\begin{figurebox}
\begin{figure}[H]
\centering
\begin{tikzpicture}
\begin{axis}[
  width=14cm, height=7cm,
  ybar stacked, bar width=18pt,
  xlabel={\rl{گروه}},
  ylabel={\rl{درصد}},
  ymin=0, ymax=105,
  symbolic x coords={%
    {\rl{شیعه‌ی عرب}},
    {\rl{سنی عرب}},
    {\rl{کردها}},
    {\rl{میانگین کل}}},
  xtick=data,
  xticklabel style={font=\small},
  legend pos=outer north east,
  legend style={font=\small, cells={anchor=west}},
  title={\rl{\textbf{«ابتدا خود را چه می‌دانید؟» (\lr{Zogby} ۲۰۰۸)}}},
  title style={font=\small},
  nodes near coords,
  nodes near coords style={font=\tiny},
]
\addplot[fill=green!50] coordinates {
  ({\rl{شیعه‌ی عرب}},40)
  ({\rl{سنی عرب}},30)
  ({\rl{کردها}},15)
  ({\rl{میانگین کل}},32)
};
\addlegendentry{\rl{ابتدا عراقی}}
\addplot[fill=red!50] coordinates {
  ({\rl{شیعه‌ی عرب}},35)
  ({\rl{سنی عرب}},45)
  ({\rl{کردها}},70)
  ({\rl{میانگین کل}},45)
};
\addlegendentry{\rl{ابتدا شیعه/سنی/کرد}}
\addplot[fill=blue!40] coordinates {
  ({\rl{شیعه‌ی عرب}},20)
  ({\rl{سنی عرب}},20)
  ({\rl{کردها}},10)
  ({\rl{میانگین کل}},18)
};
\addlegendentry{\rl{ابتدا مسلمان}}
\addplot[fill=gray!30] coordinates {
  ({\rl{شیعه‌ی عرب}},5)
  ({\rl{سنی عرب}},5)
  ({\rl{کردها}},5)
  ({\rl{میانگین کل}},5)
};
\addlegendentry{\rl{سایر}}
\end{axis}
\end{tikzpicture}
\caption{هویت فرقه‌ای/قومی غالب ─ به‌ویژه کردها (۷۰٪)}
\label{fig:stacked-identity}
\end{figure}
\end{figurebox}

\infobox{boxgreen}{نکته‌ی مثبت: رشد ناسیونالیسم عراقی}{%
هویت «عراقی» از ۲۸٪ (۲۰۰۷) به ۴۱٪ (۲۰۰۹) رشد کرد.
این روند بعدها در جنبش تشرین ۲۰۱۹ با شعار
\textbf{«نرید وطن»} (\textit{وطن می‌خواهیم}) به اوج رسید ─
نشانه‌ای از تقویت تدریجی ناسیونالیسم فرافرقه‌ای.
}


\subsection{همزیستی و ساختار حکومت: \lr{NDI}}

\begin{surveycard}{\lr{National Democratic Institute (NDI)} ─ ۲۰۰۴–۲۰۰۹}
\begin{tabularx}{\textwidth}{>{\bfseries\small}r >{\small}X}
مجری   & \lr{Greenberg Quinlan Rosner Research} \\
نمونه  & \ptho{۲}{۰۰۰}–\ptho{۲}{۴۰۰} · ۸ موج \\
هدف    & ارزیابی انتخابات و مشارکت مدنی \\
\end{tabularx}
\end{surveycard}

\begin{comparecard}[title={\textbf{\lr{NDI}: همزیستی و ساختار حکومتی مطلوب}}]
% ── ستون چپ: همزیستی ──
\begin{center}\textbf{«آیا شیعه، سنی و کرد می‌توانند\\مسالمت‌آمیز زندگی کنند؟»}\end{center}
\small
\begin{tabularx}{0.95\linewidth}{>{\bfseries\scriptsize}r *{4}{>{\centering\arraybackslash\scriptsize}X}}
\toprule
& \textbf{۰۴} & \textbf{۰۶} & \textbf{۰۸} & \textbf{۰۹} \\
\midrule
بله (کل) & ۷۵٪ & ۵۸٪ & ۵۵٪ & ۶۵٪ \\
شیعه     & ۸۰٪ & ۶۰٪ & ۵۸٪ & ۶۸٪ \\
سنی      & ۷۰٪ & ۵۲٪ & ۴۸٪ & ۵۸٪ \\
کرد      & ۶۵٪ & ۵۰٪ & ۵۵٪ & ۶۲٪ \\
\bottomrule
\end{tabularx}
\tcblower
% ── ستون راست: ساختار ──
\begin{center}\textbf{ساختار حکومتی مطلوب}\end{center}
\small
\begin{tabularx}{0.95\linewidth}{>{\bfseries\scriptsize}r *{4}{>{\centering\arraybackslash\scriptsize}X}}
\toprule
& \textbf{۰۴} & \textbf{۰۶} & \textbf{۰۸} & \textbf{۰۹} \\
\midrule
مرکزی قوی   & ۶۰٪ & ۵۵٪ & ۴۸٪ & ۵۲٪ \\
فدرالی       & ۲۵٪ & ۲۸٪ & ۳۲٪ & ۳۰٪ \\
تجزیه        & ۱۰٪ & ۱۲٪ & ۱۵٪ & ۱۲٪ \\
\midrule
\multicolumn{5}{r}{\textbf{«مرکزی قوی»}} \\
سنی  & ۸۵٪ & ۸۰٪ & ۷۵٪ & ۷۸٪ \\
شیعه & ۵۵٪ & ۵۲٪ & ۴۵٪ & ۵۰٪ \\
کرد  & ۲۰٪ & ۱۵٪ & ۱۲٪ & ۱۵٪ \\
\bottomrule
\end{tabularx}
\end{comparecard}

\infobox{boxyellow}{تحلیل: شکاف ساختاری عمیق}{%
\begin{itemize}[rightmargin=0.5cm]
\item \textbf{سنی‌ها:} ۷۵–۸۵٪ خواهان دولت مرکزی ─ ترس از تهمیش در فدرالیسم
\item \textbf{کردها:} ۳۵–۴۵٪ خواهان جدایی ─ تجربه‌ی موفق خودمختاری
\item \textbf{شیعیان:} در میانه ─ تمایل به مرکزیت با حفظ اکثریت سیاسی
\item این شکاف ریشه‌ی بسیاری از بن‌بست‌های قانون اساسی عراق است
\end{itemize}
}


\subsection{ارزش‌های فرهنگی: \lr{WVS}}

\begin{surveycard}{\lr{World Values Survey (WVS)} ─ عراق ۲۰۰۴ و ۲۰۱۳}
\begin{tabularx}{\textwidth}{>{\bfseries\small}r >{\small}X}
مدیریت  & دانشگاه‌های استکهلم و میشیگان \\
نمونه   & \ptho{۲}{۳۲۵} (۲۰۰۴) · \ptho{۱}{۲۰۰} (۲۰۱۳) \\
ویژگی   & عمیق‌ترین نظرسنجی ارزشی ─ ۲۵۰+ سؤال \\
\end{tabularx}
\end{surveycard}

\begin{table}[H]
\centering
\caption{شاخص‌های ارزشی کلیدی \lr{WVS} ─ عراق}
\label{tab:wvs-values}
\small
\begin{tabularx}{\textwidth}{>{\bfseries\small}r >{\small}X >{\centering\arraybackslash\small}p{1.3cm} >{\centering\arraybackslash\small}p{1.3cm

<!-- POSSIBLE OVERLAP DETECTED (similarity: 85%) - REVIEW BELOW -->